% ===== 原文 PDF pp.76--78（印刷 pp.75--77）=====
% [边界] 起点裁定（TIMELINE #7，Lead 已核）：p.76 首句「The one with regular puncture on a torus
%   makes manifest …」起句于 p.75 页末（"…are complementary. The one"），按「开始者写完」归 tr-06 收尾；
%   本文件自 p.76 第 3 行「Meanwhile, the construction with twisted irregular and regular punctures …」起译。
% [原文 p.75]

% [Tutor Note: p.76 开头这一段是 §4.8.3 的收尾。它对照同一个 $\cN=4$ SU(2) SYM 的两种 Class S 构造：
%   (i) $A_1$ 的 6d $(2,0)$ 理论放在带一个正则穿刺点的环面上；(ii) $A_2$ 的 6d $(2,0)$ 理论经 $\Z_2$
%   外自同构扭曲后放在球面上、带一对（扭曲正则 + 非正则）穿刺点。前者把共形流形等同于一次穿刺环面
%   的复结构模空间，后者把 SU(2) 味对称性的精细结构（含整体反常）暴露出来。]

同时，带有扭曲非正则穿刺点与扭曲正则穿刺点的构造，则使 SU(2) 味对称性的精细结构变得一目了然。特别地，$A_{2N}$ 理论中的 $\Z_2$ 扭曲正则完全穿刺点总是携带 $\mathrm{USp}(2N)$ 味对称性的 Witten 整体反常 [213]，这一结论立刻适用于实现 $\cN=4$ SYM 的 $A_2$ 情形。此外，扭曲正则穿刺点（与球面上的非正则穿刺点配对）还与 W-代数型的手征代数存在自然对应 [56]；特别地，这意味着与 $\cN=4$ SU(2) SYM 相对应的手征代数正是 $\widehat{\fsu}(2)_{-3/2}$ Kac--Moody 代数——由于该理论本身是拉格朗日式的，这一事实很容易直接验证。

\SEC{5}{IIB 型弦论中的几何工程}

上一节我们介绍了 4d $\cN=2$ SCFT 的 Class S 构造，并看到：在 Class S 构造中，Seiberg--Witten（SW）几何会自然地作为 Hitchin 可积系统的谱曲线出现。更一般地，我们在第 3 节已经看到，一个一般 4d $\cN=2$ SCFT 的库仑分支由特殊 Kähler 几何所描述。而这类几何结构也会自然地出现在 Calabi--Yau 三重复形 [262] 的模空间中，这就促使我们追问：能否直接从 II 型弦论在 Calabi--Yau 流形上的紧化出发来实现一个 4d $\cN=2$ SCFT？本节的主要目的，就是尝试对从 IIB 型弦论紧化出发的 4d $\cN=2$ SCFT 几何工程作一概述。

为了找准方向，我们首先在第 5.1 节回顾 Calabi--Yau 三重复形中的孤立三重复形超曲面奇点。接下来，我们在第 5.2 节列出由所谓 $(G,G')$ 型奇点得到的 4d $\cN=2$ SCFT 的一些典型例子。我们在第 5.3 节简要评论这一方法与 Class S 构造之间的关系。我们在第 5.4 节讨论如何把这些孤立超曲面奇点推广到其他类型的奇点。最后我们转向 F 理论，考虑 D3 膜在椭圆纤维化 Calabi--Yau 流形奇点附近的探测，并回顾 $\cN=2$ S-fold 构造方面的近期进展。

% [Tutor Note: 本节末句预告的 F 理论推广，逻辑是：IIB 的轴子--伸缩子 $\tau$ 在 F 理论中被几何化为
%   椭圆纤维的模，于是 7-膜的位置与单值性都成了纯几何数据，可以把 $\cN=2$ 的 S-fold 构造写进一个
%   椭圆纤维化的 CY 四重复形（属 §5.4.2，本块不涉及）。]

% [原文 p.76]
\SUB{5.1}{从三重复形奇点得到 4d $\cN=2$ SCFT}

% [Tutor Note: 为什么是「$M_P\to\infty$」而不是「把 CY 取紧」？因为 4d Planck 能标满足
%   $M_P^2\propto \mathrm{Vol}(W)$（$W$ 的体积越大，4d 引力越弱）；原文说「$M_P$ 正比于体积」是宽松说法，
%   严格说是 $M_P^2\propto\mathrm{Vol}(W)$。令 $M_P\to\infty$ 等价于 $\mathrm{Vol}(W)\to\infty$，即取非紧 CY。
%   这与 §3 的图像一致：只有在非紧 CY 上才留得住一个不随尺度跑的 4d 场论极限。]

为了在 IIB 型弦论于 Calabi--Yau 流形上的紧化中产生一个 4d 场论极限，通常需要把 4d Planck 能标 $M_P$ 送到无穷大（即 $M_P\to\infty$）以关掉引力相互作用；由 KK 约化可知，$M_P$ 正比于 Calabi--Yau 流形 $W$ 的体积。因此，人们应当考虑非紧 Calabi--Yau 流形上的紧化。此外，为了产生有趣的、相互作用的场论，通常还需要带奇点的 Calabi--Yau 流形。

具体起见，让我们先考虑一个带孤立典范（canonical）三重复形超曲面奇点（isolated hypersurface singularity，IHS）的 CY 流形 $W$\footnote{若簇 $X$ 的一个奇点满足以下条件，则称之为\textbf{典范奇点}（canonical singularity）：\begin{itemize}[leftmargin=1.2em,itemsep=1pt,topsep=2pt]\item 典范除子（作为 Weil 除子）$K_X$ 是 $\Q$-Cartier 的，即 $rK_X$ 是 Cartier 除子；这里 $r\in\Z$ 称为该奇点的指数（index）。\item 对任意奇点消解 $f:Y\to X$，新典范除子 $K_Y=f^*K_X+\sum_i a_iE_i$ 中的有理系数 $a_i$ 全部非负，其中 $E_i$ 表示 $Y$ 中的例外除子。\end{itemize}特别地，若所有 $a_i$ 皆为零，则该奇点容许一个 crepant 消解（crepant resolution）。若所有 $a_i$ 皆为正，这种奇点称为终端奇点（terminal singularity），近来已在 SCFT [236,263,264] 与 F 理论 [265] 中得到研究。关于这些奇点与术语的一般背景可参见例如 [266]。}，它由如下多项式 $F$ 定义：
\begin{equation}
  W := \{ F(x_1,x_2,x_3,x_4) = 0 \} \subset \C^4 \,,
  \tag{5.1}
\end{equation}
其中 $F=dF=0$ 在该孤立点处只有唯一解（不失一般性，我们设该点为原点 $x_i=0$）。注意，这类超曲面奇点中的一部分——例如下文例 5.1 中 $k$ 较大的 $A_k$ 奇点——无法嵌入紧 Calabi--Yau 之中。这也是我们一上来就取 $W$ 为非紧 CY 的另一个理由。II 型弦论在 $W$ 上紧化后的有效理论通常约化为一个 4d $\cN=2$ 规范场论\footnote{在某些情形下，为了得到一个真正的 4d $\cN=2$ 规范理论，还需要取进一步的极限以解耦多余的自由度。例如 IIA 型弦论在这样的流形上给出的是一个 4d Kaluza--Klein（KK）理论 [267]，人们需要在 Kähler 模空间中再取一个标度极限，以解耦 KK 模组成的塔、从而抵达一个 4d 规范理论；这个极限称为几何工程极限（geometric engineering limit）[27]。}。要想从 IHS 上的 IIB 型弦论得到一个 4d $\cN=2$ 超共形理论，对 $W$ 的充要条件是 [31,268]：
\begin{enumerate}[label=\arabic*.,leftmargin=2.2em,itemsep=3pt,topsep=4pt]
  \item 多项式 $F$ 必须是拟齐次的（quasi-homogeneous），即 $F$ 带有一个非平凡的 $\C^*$ 作用，使得所有坐标 $x_i$ 都具有正的权重。也就是说，该条件为
  \begin{equation}
    \C^*:\quad F(\xi^{q_i}x_i) = \xi F(x_i)\,,\qquad q_i>0\,,\qquad i=1,2,3,4\,.
    \tag{5.2}
  \end{equation}
  % [Tutor Note: (5.2) 的 $\C^*$ 作用同时给出两样东西：坐标的权重 $q_i$ 与 4d 的 $U(1)_r$ R 对称性
  %   ——取 $\xi=\ee^{\ii\theta}$ 的相位旋转即得 $U(1)_r$，而 $x_i$ 的 R 荷正是 $q_i$（归一化使 $F$ 的荷为 1）。
  %   这也是为什么「拟齐次」是出现 4d SCFT 的必要条件：没有这个 $\C^*$，就没有 $U(1)_r$，
  %   也就凑不出超共形代数里的那个 $U(1)_r$。]
  % [原文 p.77]
  这一要求来自如下事实：所考虑的 4d SCFT 具有一个 $U(1)_r$ R 对称性，它由这个 $\C^*$ 作用继承而来，因而必须被保持。
  \item 在 $\C^*$ 作用下，权重 $q_i(x_i)$ 必须满足如下条件，以保证 $W$ 上具有 Calabi--Yau 锥度量（conic metric）[31]：
  \begin{equation}
    \sum_{i=1}^{4} q_i > 1\,.
    \tag{5.3}
  \end{equation}
\end{enumerate}

第二个条件可以从描述该 IIB 型弦论背景的 2d $\cN=(2,2)$ Landau--Ginzburg（LG）模型的角度来理解：在该模型中，多项式 $F$ 被看作世界面上的超势。为看出这一点，回想：当把 IIB 型弦耦合取为 $g_s\to0$ 而保持弦标度 $\ell_s$ 固定时，IHS 处解耦的动力学由一个 4d little string theory（LST）描述（见例如 [269]），它在全息意义下对偶于背景 [270]
\begin{equation}
  \R^{1,3}\times \R_{\phi}\times \bigl(S^1\times LG(F)\bigr)/\Gamma \,,
  \tag{5.4}
\end{equation}
上的 IIB 型弦论，其中 $\Gamma$ 是恰当的 Gliozzi--Scherk--Olive（GSO）轨道折叠投影，用以保持 4d $\cN=2$ 时空超对称。这里 $\R_{\phi}$ 是一个线性伸缩子方向，其伸缩子剖面为 $\Phi:=-\tfrac{Q}{2}\phi$；$LG(F)$ 表示含四个手征超场 $x_i$、全纯超势为 $F$ 的 2d $\cN=(2,2)$ LG 模型。特别地，上述 $\C^*$ 作用被解释为该 LG 模型中的 $U(1)_r$ 对称性。进一步取低能极限 $\ell_s\to0$，便从 4d LST 重新得到一个 4d $\cN=2$ SCFT [270,271]。现在，$\cN=1$ 线性伸缩子理论的中心荷为 $\tfrac{3}{2}+3Q^2$，而 LG 模型的中心荷为 $3\hat c=3\sum_{i=1}^{4}(1-2q_i)$，剩下的 $\R^{1,3}\times S^1$ 贡献 $c=\tfrac{15}{2}$。该背景上 IIB 型弦论的自洽性要求世界面理论的总中心荷为 $26-11=15$，由此得到条件
\begin{equation}
  \frac{1}{2}Q^2 = \sum_i q_i - 1 > 0 \,,
  \tag{5.5}
\end{equation}
这与 (5.3) 一致。从几何的观点看，条件 (5.3) 保证：当该 IHS 能够嵌入紧 Calabi--Yau 三重复形时，它在 Calabi--Yau 模空间中处于有限距离处 [268]\footnote{一般地说，Calabi--Yau $d$ 重复形中的孤立超曲面奇点处于模空间有限距离处的条件是 $\hat c<d-1$ [268]。}。

% [Tutor Note: 中心荷这笔账要逐项核：(i) $\cN=1$ 线性伸缩子部分 $c=\tfrac32+3Q^2$；
%   (ii) LG 部分 $c=3\hat c=3\sum_{i=1}^4(1-2q_i)$（四个 $q_i\to0$ 的自由手征超场回到 $c=12$，
%   与每个 $(2,2)$ 手征超场 $c=2$、共四场相符）；(iii) 平坦方向 $\R^{1,3}\times S^1$ 共 5 个非紧方向，
%   在 $\cN=1$ 世界面上每个方向 = 1 个玻色子 + 1 个费米子 = $1+\tfrac12=\tfrac32$，故
%   $c=5\times\tfrac32=\tfrac{15}{2}$，这正是原文那个看似奇怪的 $15/2$ 的来源。总世界面中心荷须为
%   $26-11=15$（$bc$ 鬼 $c=-26$ 与 $\beta\gamma$ 鬼 $c=+11$ 之和为 $-15$），代入并整理即得 (5.5)。]

% [Tutor Note: 脚注 53 与 (5.3) 其实是同一个条件换了个说法：$\hat c=\sum_i(1-2q_i)=4-2\sum_i q_i$，
%   而「$d$ 重 CY 中的孤立超曲面奇点处于模空间有限距离」要求 $\hat c<d-1$；这里超曲面在 $\C^4$ 中，
%   故 $d=3$，于是 $4-2\sum_i q_i<2$，即 $\sum_i q_i>1$，正是 (5.3)。「有限距离」是物理上的关键：
%   奇点若跑到模空间无穷远，它描述的就是被解耦的 KK 塔，而不是一个 4d 场论。]
